\chapter{The specification}\label{chap:spec}

This chapter states what the main theorems of Chapter~\ref{chap:theorems} speak about: the binary
form of a coefficient list, the points of bounded height on the curve, the search region, Sturm's
theorem as a proposition and the order in which \ratpoints{} finds the points; and, from the model, the
interface of a search (its options, its refusals and what it gives back), the reference of the
delivery of the points to a callback, and the brute-force reference. The specification proper is
the four files \verb|Basic.lean|, \verb|Region.lean|, \verb|Sturm.lean| and \verb|Order.lean| in
\verb|Ratpoints/Spec|, which use Mathlib and nothing of the model, so that they can be read on their
own. The comparator challenge restates their definitions verbatim, all but the finite set of the
points (Section~\ref{sec:specBruteForce}), and imports the interface and the two references from the
model, which uses no Mathlib, as it imports the search itself; Section~\ref{sec:comparator} says
what one has to read.

The specification is the one part whose correctness no proof can check. So wherever there was a
choice, it copies what the C code computes, and each of its clauses is compared below with
\ratpoints{}~3.1.1 at the pinned commit. Where the library function \verb|find_points| and the
command-line program \verb|ratpoints| differ, as in the order of the checks on the input, the
specification and the model follow the program. The specification enters the chapters on the model
(Chapters~\ref{chap:setup} to~\ref{chap:output}) only through the statements of
Chapter~\ref{chap:theorems}.

Throughout, $c = [c_0, c_1, \dots, c_n]$ is a list of integers, the coefficient list of the input
with the constant term first, so that~$n + 1$ is its length; $d$ is its \emph{degree}, the largest
index with $c_d \ne 0$, or~$0$ if there is none; and
\[ D = d + (d \bmod 2), \qquad N = 2 \lceil n/2 \rceil \]
are~$d$ and~$n$ rounded up to even numbers. The leading zero coefficients of~$c$ \emph{lower
$\lceil n/2 \rceil$} when $\lceil d/2 \rceil < \lceil n/2 \rceil$, that is, when $D < N$, or
equivalently $d + 2 \le N$; otherwise $D = N$, or equivalently $N \le d + 1$. $H$ is a natural
number, the height bound.

\section{The binary form}\label{sec:specBinaryForm}

\begin{definition}[The binary form]\label{def:binaryForm}
  \lean{Ratpoints.degree, Ratpoints.binaryForm}\leanok
  The degree~$d$ of~$c$ is the length of~$c$ with its trailing zeros (the leading zero coefficients
  of the polynomial) dropped, less one ($0$ for a list of zeros). For integers~$x$ and~$z$,
  \[ F_c(x, z) = \sum_{i=0}^{d} c_i\, x^i z^{D-i}\,, \]
  the binary form of even degree~$D$ of the polynomial $f = \sum_i c_i X^i$.
\end{definition}

\begin{remark}[What the C code computes]\label{rem:binaryFormC}
  This is the form that the C code evaluates. The library drops the leading zero coefficients
  (\csource{find_points.c}{4624}{4632}), so that it works with the degree~$d$, and its exact check
  computes $\sum_{k \le d} c_k a^k b^{d-k}$ by Horner's rule in the numerator~$a$, on the products
  $c_k b^{d-k}$ computed once for each denominator~$b$, and multiplies by~$b$ once more when~$d$ is
  odd (\csource{find_points.c}{5291}{5298}); the result is $F_c(a, b)$. The model evaluates the form
  in the same way, in the brute-force reference (Section~\ref{sec:specBruteForce}) and in the exact
  check of the sieve (Chapter~\ref{chap:sieve}), and both evaluations are proved equal to~$F_c$.
\end{remark}

\begin{definition}[The binary form of the refusals]\label{def:binaryFormPoly}
  \uses{def:binaryForm}
  \lean{Ratpoints.binaryFormPoly, Ratpoints.eval_binaryFormPoly}\leanok
  \[ \Phi_c = \sum_{i=0}^{n} c_i X^i Z^{N-i} \in \QQ[X, Z]\,, \]
  the polynomial made homogeneous of the even degree~$N$ that the length of its coefficient list
  allows. When $\lceil n/2 \rceil \le \lceil d/2 \rceil$, the two degrees agree, $D = N$, and
  $\Phi_c(x, z) = F_c(x, z)$ for all integers~$x$ and~$z$.
\end{definition}

The refusal theorems (Theorem~\ref{thm:refusals}) speak about~$\Phi_c$ with Mathlib's
\verb|Squarefree|: no square of a non-unit divides it. The units of $\QQ[X, Z]$ are the nonzero
constants, so $4(X^2 + Z^2)$ is squarefree (\ratpoints{} accepts \verb|4 0 4|), and~$0$ is not.
$\Phi_c$ is squarefree exactly when $N \le d + 1$ and~$f$ is squarefree
(Section~\ref{sec:setupSquarefree}). The C code tests the first condition by the leading zeros
(\csource{find_points.c}{4629}{4631}) and the second in the remainder chain of the Sturm step,
which runs only when \verb|args->sturm| $\ge 0$. With the step off, \ratpoints{} refuses as not
squarefree only the forms with $d + 2 \le N$ and, when the reversal is allowed, the search region is
the whole box and~$d$ is even, those with $c_0 = c_1 = 0$ (then $X^2 \mid \Phi_c$); it accepts
other forms that are not squarefree, such as \verb|1 0 2 0 1| or \verb|0 0 1 1| under \verb|-s|.
Every input that \ratpoints{} does not refuse has $N \le d + 1$, so that there~$\Phi_c$ and~$F_c$
agree.

\section{The points}\label{sec:specPoints}

\begin{definition}[Normalized point of height at most~$H$]\label{def:IsPoint}
  \uses{def:binaryForm}
  \lean{Ratpoints.IsNormalizedPair, Ratpoints.IsPoint, Ratpoints.points, Ratpoints.mem_points, Ratpoints.finite_setOf_isPoint}\leanok
  A pair $(x, z)$ of integers is \emph{normalized and of height at most~$H$} if
  \[ \gcd(x, z) = 1, \quad z > 0 \text{ or } (z = 0 \text{ and } x = 1), \quad |x| \le H, \quad
     z \le H\,. \]
  A triple $(x, y, z)$ of integers is a \emph{normalized point of height at most~$H$} on
  $y^2 = F_c(x, z)$ if $y^2 = F_c(x, z)$ and $(x, z)$ is normalized and of height at most~$H$.
  These points form a finite set $\mathrm{points}(c, H)$: $|x| \le H$, $0 \le z \le H$, and
  $|y| \le y^2 = F_c(x, z)$.
\end{definition}

\begin{remark}[The clauses compared with \ratpoints{}]\label{rem:pointsC}
  \emph{The normalization} is the one in which the program prints the points. As~$D$ is even,
  $(-x : (-1)^{D/2} y : -z)$ is the same point as $(x : y : z)$ in the weighted projective plane
  with the weights $(1, D/2, 1)$; as both signs of~$y$ are listed, the condition can pick $z > 0$,
  and $x = 1$ when $z = 0$, where $\gcd(x, 0) = |x|$ forces $x = \pm 1$.

  \emph{The points at infinity} are therefore the triples $(1, y, 0)$ with $y^2 = c_D$: the point
  $(1 : 0 : 0)$ for odd~$d$, and $(1 : \pm\sqrt{c_d} : 0)$ for even~$d$ when~$c_d$ is a square.
  These are the points at infinity that the C code delivers (\csource{find_points.c}{4894}{4899},
  \csource{find_points.c}{5068}{5105}) with the exact check on, without \verb|RATPOINTS_NO_Y| and
  without a reversal; under \verb|RATPOINTS_NO_CHECK| it delivers one triple with $y = 0$, under
  \verb|RATPOINTS_NO_Y| only the one with~$\sqrt{c_d}$, and after a reversal the points at infinity
  of the polynomial searched, as the points $(0 : y : 1)$ of the caller's curve
  (\csource{find_points.c}{5079}{5079}; Section~\ref{sec:pointsAtInfinity}).

  \emph{Both signs of~$y$} are listed, as the program prints both; the option of the library that
  keeps only the point with $y \ge 0$ is an option of the search (Definition~\ref{def:Options}), not
  a part of the predicate. \emph{The box is closed}: $|x| = H$ and $z = H$ are included, as in the C
  code, whose numerators run over $-H, \dots, H$ and whose denominators over $1, \dots, H$.
  \emph{Coprimality} is \verb|Int.gcd|, with $\gcd(0, z) = z$, so that $x = 0$ occurs only as
  $(0, y, 1)$, as in the test \verb|relprime| of the C code (\csource{sift.c}{318}{343}), which for
  the numerator~$0$ returns whether the denominator is~$1$. \emph{The height bound applies to the
  points at infinity too} ($|1| \le H$), so that for $H = 0$ there are no points at all; the program
  refuses $H = 0$, so this never shows, and it keeps the predicate uniform.
\end{remark}

The pair alone is what the program prints when the exact check is off, as triples $(x, 0, z)$ that
need not be points; the statement on that output concludes \lname{Ratpoints.IsNormalizedPair} for
each of them (Theorem~\ref{thm:noCheck}). The predicates are decidable, so that \verb|Basic.lean|
checks facts taken from the output of the program by \verb|decide|: for instance that
$(129, 3830, 100)$, which \verb|ratpoints '-2 0 0 1' 130 -q| prints, is a point of height at
most~$130$ of $y^2 = x^3 - 2$.

\section{The search region}\label{sec:specRegion}

The caller of \verb|find_points| may restrict the search to \emph{search intervals} $[l, u]$
for~$x/z$, closed intervals whose ends are C \verb|double|s, and to a range of the denominators
from~$b_{\mathrm{low}}$ to~$b_{\mathrm{high}}$. The specification reads a double by its~$64$ bits
only to decide whether it is a NaN and to give its value in~$\QQ$ extended by~$\bot$ and~$\top$
(Mathlib's \verb|WithBotTop ℚ|); every other predicate reads the values.

\begin{definition}[The search region]\label{def:regionSpec}
  \lean{Ratpoints.IsNaN, Ratpoints.doubleValue, Ratpoints.AreSearchIntervals, Ratpoints.IntervalsWF, Ratpoints.InIntervals, Ratpoints.InRegion, Ratpoints.IsFullBox, Ratpoints.denominatorBound}\leanok
  For the bits $w < 2^{64}$ of an IEEE binary64 double, with the exponent field
  $e = \lfloor w / 2^{52} \rfloor \bmod 2^{11}$, the fraction $t = w \bmod 2^{52}$ and the sign
  $\sigma = -1$ for $w \ge 2^{63}$ and $\sigma = 1$ otherwise: $w$ is a \emph{NaN} when $e = 2047$
  and $t \ne 0$, and the \emph{value} of the double is
  \[ \mathrm{doubleValue}(w) =
    \begin{cases}
      \sigma t\, 2^{-1074} & \text{if } e = 0 \text{ (the zeros and the subnormals)}, \\
      \sigma (2^{52} + t)\, 2^{e - 1075} & \text{if } 1 \le e \le 2046, \\
      \bot \text{ for } \sigma = -1, \ \top \text{ for } \sigma = 1 & \text{if } e = 2047
    \end{cases}\,,
  \]
  so that both zeros have the value~$0$. The intervals of a list over the values \emph{are search
  intervals} when each $[l, u]$ has $l \le u$ and each upper end lies below the lower end of the
  next interval; a list of pairs of doubles, each given by its bits~$w$, is \emph{well formed} when
  no end is a NaN and the values of the ends are search intervals. A rational number~$r$ \emph{lies in} a list~$D$ of intervals
  over the values when $l \le r \le u$ for some $[l, u] \in D$. A triple $(x, y, z)$ lies \emph{in
  the search region} of the intervals~$D$ (over the values) and the range
  from~$b_{\mathrm{low}}$ to~$b_{\mathrm{high}}$ when
  \[
    z = 0 \ \lor\ \bigl(b_{\mathrm{low}} \le z \ \land\ (b_{\mathrm{high}} < 1 \lor
      z \le b_{\mathrm{high}}) \ \land\ (D = [\,] \lor x/z \text{ lies in } D)\bigr)\,:
  \]
  a point at infinity lies in every region, and~$x/z$ may lie anywhere when there are no intervals.
  The search region is the \emph{whole box} of the height bound~$H$ when there are no intervals,
  $b_{\mathrm{low}} \le 1$, and $b_{\mathrm{high}} < 1$ or $H \le b_{\mathrm{high}}$. The
  \emph{largest denominator} of the search is $b_{\mathrm{high}}$ when $1 \le b_{\mathrm{high}} \le
  H$, and~$H$ otherwise.
\end{definition}

The statements take the search intervals of the options, a list of pairs of doubles of the model
(Definition~\ref{def:Double}), by the values of their ends where they speak about the points and by
their bits where they test well-formedness (Definition~\ref{def:Options}). The value of a NaN is
junk, an infinity, and nothing depends on it: the statements that read the values assume a
successful search, which refuses a NaN end. \verb|Region.lean| checks its definitions by
\verb|decide +kernel|, which adds no axiom, on a few doubles and on intervals that \ratpoints{} accepts
and refuses: $[-0.0, 0.0]$ and $[0.0, -0.0]$ are accepted, and $[-1, -0.0]$ followed by $[0.0, 1]$
is refused.

\begin{remark}[The region compared with \ratpoints{}]\label{rem:regionC}
  The ends of the intervals are C doubles (\verb|ratpoints_interval|,
  \csource{ratpoints.h}{359}{359}). The C code does no floating-point arithmetic on them that could
  decide a point: it compares them, converts between \verb|long|s and doubles, takes doubles apart
  and scales them by powers of~$2$ (Chapter~\ref{chap:region}), and computes with them only in the
  estimate of the work in the choice of the moduli (Chapter~\ref{chap:moduli}). So the
  specification can speak about their exact values.

  Well-formedness is the test of \verb|find_points_work_1| that refuses intervals as bad arguments
  (\csource{find_points.c}{4851}{4854}): an interval without $l \le u$, or one that does not begin
  above the end of the interval before it. Every comparison with a NaN is false, so that a NaN end
  is refused too; infinite ends, degenerate intervals $[a, a]$ and any number of intervals are
  accepted. On doubles that are not NaNs the comparisons of the C language agree with the order of
  the values, with~$-0.0$ equal to~$0.0$; so the region is the same for~$-0.0$ and~$0.0$, for the
  specification as for the C code.

  Without intervals the library searches $[-h, h]$ for the double~$h$ of~$H$, which holds
  every~$x/z$ of the box (\csource{find_points.c}{4831}{4841}). It takes $b_{\mathrm{low}} < 1$
  as~$1$ and a $b_{\mathrm{high}}$ outside $1, \dots, H$ as~$H$
  (\csource{find_points.c}{4659}{4661}), which the region states without clamping; it delivers the
  points at infinity whatever the region, and forbids the reversal exactly when the region is not
  the whole box (\csource{find_points.c}{4686}{4687}; Chapter~\ref{chap:region}).
\end{remark}

\begin{definition}[Sturm's theorem]\label{def:sturmTheorem}
  \lean{Ratpoints.SturmTheorem}\leanok
  The proposition that for every squarefree real polynomial~$p$ and all reals $a \le b$
  \[
    \#\{t \in \mathbb{R} : a < t \le b,\ p(t) = 0\} + V_p(b) = V_p(a)\,,
  \]
  where $V_p(t)$ is the number of sign changes (\verb|List.signVariations|) of the values at~$t$ of
  the Sturm sequence of~$p$ and its derivative (\verb|Polynomial.sturmSeq p (derivative p)|), and
  the number of roots is \verb|Set.ncard|.
\end{definition}

The Sturm step of \ratpoints{} narrows the search intervals to those on which~$f$ may take a value
$\ge 0$ (Chapter~\ref{chap:region}); that it keeps every point rests on Sturm's theorem, which
Mathlib does not have, though it has the Sturm sequence and the sign changes of a list. The form for
$(a, b]$ is the one that the bisection of the step needs, as it evaluates the sequence at midpoints
that may be roots of~$f$. The three certified statements that need it (membership, the order and the
coverage when the exact check is off) take the hypothesis $0 \le \mathrm{sturm} \to
\mathrm{SturmTheorem}$, so that they hold outright with the Sturm step off. It is a hypothesis, not
an axiom; a proof of the theorem would discharge it without changing a statement.

\section{The order of \ratpoints{}}\label{sec:specOrder}

The file \verb|Order.lean| defines the list of the points in the order in which \ratpoints{} finds
them, in terms of the polynomial, and the order theorem (Theorem~\ref{thm:order}) says that a search
with the exact check on returns this list, filtered by~$\mathrm{noY}$ and the search region,
assuming Sturm's theorem when the Sturm step is on. It uses the prime factors of a number
(\verb|Nat.primeFactorsList|, in increasing order with multiplicity).

\begin{definition}[The points in the order of \ratpoints{}]\label{def:orderedPoints}
  \uses{def:binaryForm, def:IsPoint}
  \lean{Ratpoints.searchedPoly, Ratpoints.FactorsByTrialDivision, Ratpoints.binaryDivisors, Ratpoints.searchDenominators, Ratpoints.pointsWithDen, Ratpoints.orderedPoints}\leanok
  For the switch~$r$ of the reversal and the coefficient list~$c$ of degree~$d$, the polynomial that
  \ratpoints{} searches and whether it is the reversed one, $\mathrm{searchedPoly}(r, c)$, is the first
  of the following that applies:
  \begin{enumerate}
    \item $((c_d, c_{d-1}, \dots, c_0), \mathrm{true})$ if~$r$ is true, $d$ is even and $c_0 = 0$;
    \item $((c_d, c_{d-1}, \dots, c_0), \mathrm{true})$ if~$r$ is true, $d$ is even, $c_d$ is a
      square and~$c_0$ is not a square;
    \item $((0, c_d, c_{d-1}, \dots, c_1), \mathrm{true})$ if~$r$ is true, $d$ is odd,
      $|c_d| \ne 1$, $c_0 = 0$ and $|c_1| = 1$;
    \item $(c, \mathrm{false})$ otherwise.
  \end{enumerate}
  A natural number~$n$ is \emph{factored by trial division} when the product of its prime factors
  from~$1021$ on, with multiplicity, is below $1021^2$, or none of its prime factors exceeds~$1021$.
  The squarefree divisors of~$n$ at most~$B$ \emph{in the order of binary counting}: for the
  distinct prime factors $p_1 < p_2 < \cdots$ of~$n$, the products of the sublists of
  $(p_1, p_2, \dots)$ in the order of \verb|List.sublists|,
  $1, p_1, p_2, p_1 p_2, p_3, p_1 p_3, \dots$, with~$p_1$ the lowest bit, those at most~$B$. The
  \emph{denominators of the search} of~$c$ up to~$B$, for a bound~$m$ of the number of divisors:
  when~$d$ is odd, $|c_d| \ne 1$, $|c_d|$ is factored by trial division and has fewer than~$m$
  squarefree divisors at most~$B$, the numbers $e k^2 \le B$, divisor~$e$ by divisor in the order of
  binary counting and $k = 1, 2, \dots$ for each; otherwise $1, \dots, B$. The \emph{points with the
  denominator}~$z$: for $x = -H, \dots, H$ in turn, with $y = \lfloor \sqrt{F_c(x, z)} \rfloor$ ($0$
  when $F_c(x, z) < 0$), those among $(x, y, z)$ and $(x, -y, z)$ that are normalized points of
  height at most~$H$, in this order, one triple when $y = 0$; the denominator~$0$ gives the points
  at infinity. The \emph{points in the order of \ratpoints{}}, $\mathrm{orderedPoints}(m, r, c, H, B)$:
  for $(c', \mathrm{rev}) = \mathrm{searchedPoly}(r, c)$, the points of~$c'$ with the
  denominator~$0$ and then with the denominators of the search of~$c'$ up to~$B$, in turn, each
  $(a, y, b)$ mapped to $(b, y, a)$ for $a \ge 0$ and to $(-b, y, -a)$ for $a < 0$
  when~$\mathrm{rev}$ is set.
\end{definition}

The binary form of a reversed list is that of~$c$ with~$x$ and~$z$ exchanged, except for even~$d$
with $c_0 = c_1 = 0$, which \ratpoints{} refuses; so the map at the end is a bijection from the points
of~$c'$ to those of~$c$ that preserves the box (Chapter~\ref{chap:setup}). In the order theorem,
$m$ is $\mathrm{maxDivisors}$ (Definition~\ref{def:maxDegree}), $r$ the switch of the reversal when
the search region is the whole box and false otherwise, and~$B$ the largest denominator: the search
intervals and the least denominator only filter the list.

\begin{remark}[The order compared with \ratpoints{}]\label{rem:orderC}
  $\mathrm{searchedPoly}$ is the reversal check of \verb|find_points_work_1|
  (\csource{find_points.c}{4732}{4803}) on the coefficient list; items (1) to~(3) are its cases~1
  to~3. In case~(1) the check refuses the form as not squarefree when also $c_1 = 0$, since then
  $X^2 \mid \Phi_c$ (\csource{find_points.c}{4744}{4749}), and otherwise reverses $c_0, \dots, c_d$,
  which lowers the degree to~$d - 1$; in case~(2) it reverses $c_0, \dots, c_d$, in case~(3)
  $c_1, \dots, c_d$. The squares are those of GMP's \verb|mpz_perfect_square_p|, which agrees with
  Mathlib's \verb|IsSquare| on~$\ZZ$: no negative integer is a square. After a reversal the C code
  hands each point $(a, y, b)$ of the polynomial searched to the caller mapped as above
  (\csource{find_points.c}{5312}{5316}).

  The C code delivers the points at infinity first, then sieves the denominators in the order of its
  loop, and for each denominator delivers the points by increasing numerator, $(a, y, b)$ before
  $(a, -y, b)$ (\csource{find_points.c}{5300}{5327}). For odd~$d$ the loop departs from
  $1, \dots, B$ in two ways. When $|c_d| = 1$ it visits only the squares, in increasing order; the
  specification lists $1, \dots, B$, and the denominators that are not squares carry no point. When
  $|c_d| \ne 1$, \verb|setup_us1| factors~$|c_d|$ by trial division up to~$1021$, which succeeds
  exactly when~$|c_d|$ is factored by trial division in the sense above, and makes the squarefree
  divisors up to~$b_{\mathrm{high}}$ in the order of binary counting
  (\csource{find_points.c}{777}{860}, \csource{find_points.c}{864}{881}); when there are fewer than
  \verb|MAX_DIVISORS|, the loop visits~$e k^2$, divisor by divisor and~$k$ increasing
  (\csource{find_points.c}{3937}{3940}), the one case in which the denominators do not increase.
  Chapter~\ref{chap:denominators} proves that the loop of the model visits the denominators in this
  order.
\end{remark}

\verb|Order.lean| checks its definitions by \verb|decide +kernel| on outputs of the C program, among
them \verb|ratpoints '-1 0 9' 10 -q|, which reverses the polynomial to \verb|9 0 -1|, and
\verb|ratpoints '-6 5 -3 30' 12 -q|, which visits the denominator~$6$ before~$5$, in the order
$1, 2, 3, 6, 5, 10, 15, 30$ of the divisors of~$30$, and on runs with search regions.

\section{The options, the refusals and the outcome of a search}\label{sec:specInterface}

The main theorems speak about the search $\mathrm{findPoints}(c, H, s, o)$, which returns the list
of the points or a refusal, and about the search with a callback (Definition~\ref{def:findPoints}).
Their arguments and results are types of the model, whose fields mirror the fields of the arguments
\verb|ratpoints_args| of the library (\csource{ratpoints.h}{362}{400}) and the bits of its flags
(\csource{ratpoints.h}{402}{424}).

\begin{definition}[The options, the refusals and the outcome of a search]\label{def:Options}
  \uses{def:Double, def:regionSpec}
  \lean{Ratpoints.Options, Ratpoints.SieveOptions, Ratpoints.Refusal, Ratpoints.Outcome, Ratpoints.defaultSturm, Ratpoints.domainValues, Ratpoints.domainBits}\leanok
  The \emph{options} of a search are those of the library that change the refusals, the order of the
  output or what is delivered, with the defaults of a run of the program without options: whether
  the set-up may reverse the polynomial, that is, whether the bit \verb|RATPOINTS_NO_REVERSE| is
  clear (true by default); the depth \verb|args->sturm| of the Sturm step, an integer, negative for
  none, by default $\mathrm{defaultSturm} = 10$ (\verb|RATPOINTS_DEFAULT_STURM|); the two options of
  the output $\mathrm{noCheck}$ and~$\mathrm{noY}$, false by default, the bits
  \verb|RATPOINTS_NO_CHECK|, under which every numerator that passes the sieve is delivered with
  $y = 0$ without the exact check, and \verb|RATPOINTS_NO_Y|, under which only the point with
  $y \ge 0$ of a pair is delivered; and the search region: the search intervals, a list of pairs of
  doubles for the array \verb|domain| with its length \verb|num_inter|, empty by default, and the
  least and the largest denominator, integers for the \verb|long|s \verb|b_low| and \verb|b_high|,
  by default~$1$ and~$0$ (\csource{main.c}{348}{348}, \csource{main.c}{352}{352},
  \csource{ratpoints.h}{403}{405}); the program sets \verb|b_high| to~$H$, and the library takes a
  value below~$1$ as~$H$, so that the two defaults give the same search. The statements read the
  search intervals by the values of their ends, $\mathrm{domainValues}$, and by their bits,
  $\mathrm{domainBits}$.

  The \emph{options of the sieve} are the values of the options of the program that tune the sieve,
  which change neither the points nor their order: the number of primes examined (\verb|-p|), the
  cut points of the stages (\verb|-n|, \verb|-N|), the parameters of the choice of the moduli and of
  its correction (\verb|-r -R -U -C -A -P -Q -W|), the number of forbidden divisors (\verb|-F|), the
  array size, and whether the Jacobi test runs (\verb|-j|, \verb|-J|), as the program reads them
  into \verb|ratpoints_args|, $-1$ for an absent option that takes a number
  (\csource{main.c}{337}{347}, \csource{main.c}{349}{350}, \csource{ratpoints.h}{365}{367}).

  The \emph{refusals} are five, each with the error of the program: more than
  $\mathrm{maxDegree} + 1$ coefficients (error~3), fewer than two (error~5), the height bound~$0$ or
  one of~$2^{63}$ or more, which is not a C \verb|long| (error~4), a binary form that is not
  squarefree (error~8), and search intervals that are not well formed (error~9).

  The \emph{outcome} of a search that is not refused is what the library gives back to its caller:
  the data of the callback, the count, the three bits of the flags that the program reads after the
  search, \verb|RATPOINTS_REVERSED|, \verb|RATPOINTS_USE_SQUARES| and \verb|RATPOINTS_USE_SQUARES1|,
  the search intervals as the search leaves them, and the numbers of the moduli used in the stages
  (\csource{ratpoints.h}{377}{383}, \csource{ratpoints.h}{414}{414},
  \csource{ratpoints.h}{416}{417}).
\end{definition}

The search refuses in the order of the program, which differs from that of the library on a few
inputs; Section~\ref{sec:setupOrder} gives this order with its lines of C, and
Theorem~\ref{thm:refusals} says exactly which inputs are refused, with the Sturm step on and off.
The options of the sieve are not options in the sense above: the main theorems hold for every value
of them, with no hypothesis, because the search chooses its moduli from the curve and these values
only steer the choice (Chapter~\ref{chap:moduli}). The search normalizes both kinds of options as
\verb|find_points_work_1| does (\lname{Ratpoints.Options.normalize},
\lname{Ratpoints.SieveOptions.normalize}), and the statements speak about the options as the caller
gives them, as \verb|find_points_work| restores the fields that it normalized.

\begin{definition}[The bounds that the statements mention]\label{def:maxDegree}
  \lean{Ratpoints.maxDegree, Ratpoints.maxDivisors, Ratpoints.fitsLong}\leanok
  $\mathrm{maxDegree} = 100$, the largest degree of a polynomial that \ratpoints{} accepts, the value
  of \verb|RATPOINTS_MAX_DEGREE| (\csource{ratpoints.h}{33}{33}): a list of coefficients has at most
  $\mathrm{maxDegree} + 1 = 101$ entries. $\mathrm{maxDivisors} = 512$, the value of
  \verb|MAX_DIVISORS| (\csource{find_points.c}{121}{121}): the loop over the denominators of a
  polynomial of odd degree runs over the divisors of its leading coefficient times squares only when
  there are fewer than this many divisors up to the largest denominator. An integer is a C
  \verb|long| when it lies in $[-2^{63}, 2^{63})$ (\csource{rp-private.h}{41}{43}): \ratpoints{} needs
  a \verb|long| of~$64$ bits, and the height bound is one.
\end{definition}

Like the other constants that the C code fixes when it is compiled, the two are definitions of the
model with the values of the default build, and no proof depends on their values beyond
$0 < \mathrm{maxDegree} < 2^{57}$; the author checks this by building the project with other values
of the constants (Chapter~\ref{chap:ccode}).

\section{The reference of the delivery}\label{sec:specDelivery}

The library function \verb|find_points| returns no list of points. For each point that it finds it
calls a function \verb|process| of the caller with the coordinates, the caller's data \verb|info|
and a pointer to a flag \verb|quit|; it adds the value that the call returns to a count
\verb|total|, which it returns, and it stops the search as soon as a call sets the flag
(\csource{ratpoints.h}{436}{437}).

\begin{definition}[The reference of the delivery]\label{def:deliver}
  \lean{Ratpoints.Answer, Ratpoints.Delivery, Ratpoints.Delivery.call, Ratpoints.Delivery.deliver}\leanok
  The callback is a function of~$x$, $z$, $y$ and the data, in the order of the C code, that returns
  what the call gives: the weight that it returns, an integer, whether it set \verb|*quit|, and the
  data after the call (\csource{ratpoints.h}{436}{437}, \csource{ratpoints.h}{441}{442}). A
  \emph{delivery} is the data of the callback, the count, the sum of the weights so far, and the
  flag of the stop; one call on a delivery adds the weight to the count and takes the flag and the
  data that the call leaves. The \emph{delivery of a list}~$L$ of points to the callback, from a
  delivery~$r$: for each point $(x, y, z)$ of~$L$ in turn, while the flag is clear, the call with
  $(x, z, y)$; the points after a call that sets the flag, and all of them when it is set at the
  start, are passed over.
\end{definition}

This fold over a list is a reference of the model, as the brute-force reference is, and models no C
code: the search delivers each point where it finds it and tests the flag where the C code tests
\verb|quit| (Chapter~\ref{chap:output}). By the protocol theorem (Theorem~\ref{thm:protocol}) the
search with any callback refuses what $\mathrm{findPoints}$ refuses and otherwise gives the data and
the count of the delivery of its list from the caller's data, the count~$0$ and the flag clear;
$\mathrm{findPoints}$ is itself the search with the callback that appends each point to an array,
with the weight~$1$, and never stops (\lname{Ratpoints.appendPoint}). The weights are integers
without a bound, for the \verb|int|s that the callback returns.

\section{The brute-force reference}\label{sec:specBruteForce}

\begin{definition}[The brute-force reference]\label{def:bruteForce}
  \lean{Ratpoints.bruteForce, Ratpoints.BruteForce.pointsOfDen}\leanok
  The list, computed in the model, of the triples found by testing, in this order, the candidate
  $(1, 0)$ if $H \ge 1$, and then the pairs $(x, z)$ with $1 \le z \le H$, $-H \le x \le H$ and
  $\gcd(x, z) = 1$, with~$z$ increasing and, for each~$z$, $x$ increasing. For each candidate,
  $v = F_c(x, z)$ is computed by Horner's rule in~$x$ on the products $c_k z^{d-k}$, computed once
  for each~$z$, times~$z$ when~$d$ is odd; if $v \ge 0$ and $v = y^2$ for
  $y = \lfloor \sqrt{v} \rfloor$, the triple $(x, y, z)$ is listed, followed by $(x, -y, z)$ if
  $y \ne 0$. The part of the list over one denominator~$z \ge 1$, or over the candidate $(1, 0)$
  for $z = 0$, is a function of its own, with which the proofs compare a search denominator by
  denominator.
\end{definition}

The reference is the core of the script \verb|verify-test4.py| of the C repository, written in the
model, in Lean core without Mathlib. It evaluates the form as the exact check of the C code does
(Remark~\ref{rem:binaryFormC}), so that the candidate $(1, 0)$ needs no case of its own, and it
computes the degree by the expression of $\mathrm{degree}$, which the model cannot import. It takes
no hypothesis on~$c$ or~$H$. The proofs of the main theorems compare the list of the search with
this reference, denominator by denominator (Chapter~\ref{chap:theorems}); for a polynomial that
\ratpoints{} neither reverses nor searches in the order of the divisors, and for the whole box with the
options of the output clear, its list is that of \ratpoints{}, in the same order.

\begin{theorem}[The brute-force reference finds exactly the points]\label{thm:bruteForce}
  \uses{def:bruteForce, def:IsPoint}
  \lean{Ratpoints.mem_bruteForce_iff, Ratpoints.nodup_bruteForce, Ratpoints.toFinset_bruteForce}\leanok
  For every list~$c$ and every~$H$: a triple lies in $\mathrm{bruteForce}(c, H)$ if and only if it
  is a normalized point of height at most~$H$ on $y^2 = F_c(x, z)$; the list has no duplicates; and
  its set of elements is $\mathrm{points}(c, H)$.
\end{theorem}

\begin{proof}
  \uses{def:binaryForm}
  \leanok
  Horner's rule on the products, times~$z$ for odd~$d$, is $F_c(x, z)$; $v$ is the square of~$y$
  exactly when $v \ge 0$, $v = \lfloor \sqrt{v} \rfloor^2$ and $y = \pm \lfloor \sqrt{v} \rfloor$.
  The candidate $(1, 0)$ gives the points with $z = 0$, since $\gcd(x, 0) = 1$ and the
  normalization force $x = 1$, and the pairs give those with $z > 0$. The triples of different
  candidates differ in~$x$ or in~$z$, and $y \ne -y$ for $y \ne 0$, so that there are no duplicates.
  The third statement follows from the first.
\end{proof}

Comparator certifies the first two statements. The third is not in the challenge, because
$\mathrm{points}$ is not restated there: comparator compares the full declarations that a certified
statement reaches, so that the finiteness proof inside $\mathrm{points}$ would have to be reproduced
term for term, and the first two statements say everything that the third says.
